% Generated draft: requires mathematical and source review.
\documentclass[11pt]{article}
\usepackage[margin=25mm]{geometry}
\usepackage{fontspec}
\setmainfont{Latin Modern Roman}
\usepackage{amsmath,amssymb,mathtools,unicode-math}
\setmathfont{Latin Modern Math}
\usepackage{enumitem,needspace,xurl,hyperref}
\setlength{\parindent}{0pt}
\setlength{\parskip}{5pt}
\setlength{\emergencystretch}{3em}
\newcommand{\N}{\mathbb{N}}
\newcommand{\Z}{\mathbb{Z}}
\newcommand{\Q}{\mathbb{Q}}
\newcommand{\R}{\mathbb{R}}
\newcommand{\C}{\mathbb{C}}
\newcommand{\F}{\mathbb{F}}
\newcommand{\A}{\mathbb{A}}
\newcommand{\PP}{\mathbb{P}}
\newcommand{\E}{\mathbb{E}}
\newcommand{\eps}{\varepsilon}
\newcommand{\dd}{\,\mathrm{d}}
\begin{document}
\section*{UnsolvedMath ID: 30002356}
\textbf{Source problem code:} catalogue entry
\subsection{Definitions and mathematical statement}
The supplied record asks for a general computation, but does not specify a base field, a range of degrees, or all coefficient conventions. The following defines the requested object under the ordinary Galois-cohomology interpretation; these conventions require confirmation against the source.

Let \(k\) be a field, let \(k_s\) be a separable closure, and write \(\Gamma_k=\operatorname{Gal}(k_s/k)\). Let \(G\) be a simple algebraic group over \(k\), and let \(M\) be a discrete abelian group with continuous \(\Gamma_k\)-action. For each field extension \(K/k\), let \(M_K\) denote the corresponding restricted Galois module and let \(H^n(K,M_K)\) denote continuous Galois cohomology. Here \(n\geq 0\) is a fixed cohomological degree; the source does not specify which degrees are intended. Write \(H^1(K,G)\) for the pointed set of isomorphism classes of \(G_K\)-torsors, and write \(1_K\) for the class of the trivial torsor.

A degree-\(n\) cohomological invariant with coefficients in \(M\) is a family of maps
\[
 a_K:H^1(K,G)\longrightarrow H^n(K,M_K),\qquad K/k\text{ a field extension},
\]
natural under embeddings of fields over \(k\). Thus, for every such embedding \(K\hookrightarrow L\) and every \(\xi\in H^1(K,G)\),
\[
 a_L(\xi_L)=\operatorname{res}_{L/K}\bigl(a_K(\xi)\bigr).
\]
The invariant is normalized if \(a_K(1_K)=0\) for every extension \(K/k\). Pointwise addition makes the normalized invariants an abelian group, denoted
\[
 \operatorname{Inv}^{n}(G,M)^{\mathrm{norm}}.
\]

\textbf{Requested computation.} Determine \(\operatorname{Inv}^{n}(G,M)^{\mathrm{norm}}\) for simple algebraic groups \(G\) and Galois modules \(M\), including the basic coefficient case \(M=\Z/2\Z\). This retains the generality requested in the supplied record; it does not assert a computation or a uniform classification.

\textbf{Unresolved qualifications.} The record does not specify restrictions on \(k\), whether simplicity is over \(k\) or absolute, or whether particular isogeny types are intended. It also does not specify coefficient twists or whether a cohomology theory other than ordinary Galois cohomology is intended in positive characteristic. In particular, the meaning and intended scope of the coefficient case \(\Z/2\Z\) in characteristic \(2\) require confirmation. The definitions above describe one conventional reading without resolving these omissions.
\subsection{Short English statement}
Compute the abelian group of natural rules assigning a degree-\(n\) cohomology class to each torsor under a simple algebraic group, with value zero on every trivial torsor. The request includes coefficients \(\Z/2\Z\). The base field, intended notion of simplicity, degrees, and coefficient conventions remain unspecified. The supplied record labels the problem partially solved; this draft does not independently verify that assessment.
\subsection{Sources}
\begin{itemize}
\item Local registry: \texttt{data/unsolvedmath/problems.json}, permanent UnsolvedMath ID \(30002356\), source problem code \texttt{catalogue entry}, title \emph{Cohomological Invariants of Simple Algebraic Groups}, in the supplied Oberwolfach Reports Open Problems collection. The record gives proposed year \(2013\) and supplies no named proposer.
\item Supplied source citation: \emph{Quadratic Forms and Linear Algebraic Groups (2014)}, Oberwolfach Reports, pp.~1819--1859, DOI \texttt{10.4171/owr/2013/31}. \url{https://doi.org/10.4171/owr/2013/31}. The citation's year \(2014\) is retained as supplied; the proposed year and DOI contain \(2013\), and their relationship has not been verified.
\item Supplied literature source, labelled ``Oberwolfach report on cohomological invariants'': \url{https://oa.tib.eu/renate/server/api/core/bitstreams/75e13110-0aa1-46a6-8c43-3bf00c863bc2/content}.
\item Supplied literature source, labelled ``Merkurjev, degree-three invariants'': \url{https://doi.org/10.4171/JEMS/599}.
\item Literature assessment dated \texttt{2026-08-21}: the supplied curation labels the problem partially solved and states that the general computation remains open, while degree-three cohomological invariants are classified in major families, including adjoint inner-type groups. It describes this as substantial partial progress rather than a complete computation. This is supplied dated curation derived from the OWR material, not an independently verified literature assessment; its status and cited claims require independent checking.
\item The supplied registry labels the statement \texttt{corrected\_verified} and says that notation, typography, and wording were normalized without changing mathematical content. This is a source-record qualification, not mathematical verification performed for this draft.
\end{itemize}
\end{document}
